\subsection{Grandes relations de l'optique géométrique}

L'objectif de cette partie est de présenter les relations de conjugaison 
du modèle de l'optique géométrique
qui servent de base à la démonstration des relations demandées 
dans les parties suivantes.

Dans tout ce qui suit, 
on se place dans les conditions de Gauss
et on considère une lentille mince de centre optique $O$, 
de foyers objet $F$ et image $F'$.

\subsubsection{Relation de conjugaison de Descartes}

Cette relation exprime la position de l'image $A'$ 
par rapport à la position de l'objet $A$, 
en utilisant le centre optique $O$ comme origine des mesures :

\begin{equation}
    \frac{1}{\overline{OA'}} - \frac{1}{\overline{OA}} = \frac{1}{f'}
\end{equation}

Où :
\begin{itemize}
    \item $\overline{OA}$ représente la position de l'objet ;
    \item $\overline{OA'}$ représente la position de l'image ;
    \item $f'$ est la distance focale image de la lentille.
\end{itemize}

\subsubsection{Relation de conjugaison de Newton}

Cette relation utilise les foyers $F$ et $F'$ comme origine. 
Elle est particulièrement utile pour la méthode de Badal :

\begin{equation}
    \overline{FA} \cdot \overline{F'A'} = f \cdot f' = -f'^{2}
\end{equation}

Ces relations fondamentales permettent d'établir les expressions propres 
aux différentes méthodes de focométrie expérimentale.